\documentclass[11pt]{article}
\usepackage[a4paper,margin=24mm]{geometry}
\usepackage{fontspec}
\usepackage{microtype}
\usepackage{enumitem}
\usepackage{xcolor}
\usepackage{fancyhdr}
\usepackage[hidelinks]{hyperref}
\setmainfont{TeX Gyre Pagella}
\definecolor{statusgreen}{HTML}{EAF6EE}
\pagestyle{fancy}
\fancyhf{}
\fancyfoot[C]{\parbox{0.94\textwidth}{\centering\scriptsize Hiika hojii kan hin gamaaggamamne. OpenStax, \textit{Prealgebra 2e}, \S2.1 irraa hiikame; CC BY-NC-SA 4.0.\\\href{https://openstax.org/books/prealgebra-2e/pages/2-1-use-the-language-of-algebra}{OpenStax irratti bilisaan dubbisi.}}}
\setlength{\headheight}{14pt}
\setlist[itemize]{leftmargin=1.6em,itemsep=0.35em}

\begin{document}

\begin{center}
  {\Large\bfseries Afaan Aljebiraa}\par
  \vspace{0.35em}
  {\small OpenStax \textit{Prealgebra 2e}, kutaa 2.1 --- barruu filatame}
\end{center}

\vspace{0.6em}
\colorbox{statusgreen}{\parbox{0.94\textwidth}{\small
Kun hiika hojii kan ammallee gamaaggama afaanii yookiin hawaasaa hin arganne dha.
Jechoonni as keessatti fayyadaman kaadhimamtoota ragaa maddaa irraa baafaman;
hin raggaasifamne.}}

\section*{2.1 Afaan aljebiraa fayyadamuu}

Dhuma kutaa kanaatti, wantoota armaan gadii gochuu ni dandeessa:

\begin{itemize}
  \item jijjiiramootaa fi mallattoolee aljebiraa fayyadamuu;
  \item ibsamootaa fi himoota walqixaa adda baasuu.
\end{itemize}

Aljebiraa keessatti qubeewwan jijjiiramoota bakka buusuuf fayyadu. Qubeewwan
$x$, $y$, $a$, $b$, fi $c$ yeroo baay'ee jijjiiramoota bakka buusuuf fayyadu.

\subsection*{Jijjiiramaa fi dhaabataa}

\textbf{Jijjiiramaan} qubee lakkoofsa yookiin hamma gatiin isaa jijjiiramuu
danda'u bakka bu'u dha.

\textbf{Dhaabataan} lakkoofsa gatiin isaa yeroo hunda waluma fakkaatee turu dha.

\subsection*{Ibsamootaa fi himoota walqixaa}

\textbf{Ibsamni} lakkoofsa, jijjiiramaa, yookiin walitti qabama lakkoofsotaa,
jijjiiramootaa fi mallattoolee qoyyabaa ti.

\textbf{Himni walqixaa} ibsamoota lama mallattoo walqixaatiin walitti hidhaman
irraa ijaarama.

\vfill
{\footnotesize\textbf{Daangaa haalaa:} Hiikni kun raggaasisa jechootaa,
gamaaggama nama afaan dhalootaa dubbatuu, yookiin mirkaneessa hawaasaa miti.}

\end{document}
